\documentclass[11pt]{article}

%%%%%%%%%%%%%%%%%%%%%%%%%%%%%%%%%%%%%%%%
% PACKAGES
%%%%%%%%%%%%%%%%%%%%%%%%%%%%%%%%%%%%%%%%
\usepackage[margin=1in]{geometry}
\usepackage{amsmath,amsfonts,amssymb}
\usepackage{graphicx}
\usepackage{xcolor}
\usepackage{tikz}
\usetikzlibrary{arrows.meta, positioning, shapes.geometric, calc}
\usepackage{hyperref}
\usepackage{enumitem}
\usepackage{titlesec}
\usepackage{array}
\usepackage{booktabs}

%%%%%%%%%%%%%%%%%%%%%%%%%%%%%%%%%%%%%%%%
% COLOR SCHEME (MEMORY OPTIMIZED: black/blue/red)
%%%%%%%%%%%%%%%%%%%%%%%%%%%%%%%%%%%%%%%%
\definecolor{myblue}{RGB}{20,90,160}
\definecolor{myred}{RGB}{180,40,40}
\newcommand{\blue}[1]{\textcolor{myblue}{#1}}
\newcommand{\red}[1]{\textcolor{myred}{#1}}

\hypersetup{
  colorlinks=true,
  linkcolor=myblue,
  urlcolor=myblue,
  citecolor=myblue
}

%%%%%%%%%%%%%%%%%%%%%%%%%%%%%%%%%%%%%%%%
% SECTION STYLING
%%%%%%%%%%%%%%%%%%%%%%%%%%%%%%%%%%%%%%%%
\titleformat{\section}
{\large\bfseries\color{myblue}}
{\thesection}{1em}{}

\titleformat{\subsection}
{\bfseries\color{myblue}}
{\thesubsection}{1em}{}

%%%%%%%%%%%%%%%%%%%%%%%%%%%%%%%%%%%%%%%%
% MEMORY BOXES
%%%%%%%%%%%%%%%%%%%%%%%%%%%%%%%%%%%%%%%%
\newenvironment{bluebox}
{\vspace{6pt}\noindent\begin{tikzpicture}
\node[
draw=myblue, line width=1pt,
rounded corners=6pt,
inner sep=8pt,
text width=\dimexpr\linewidth-20pt\relax
]\bgroup
\begin{minipage}{\dimexpr\linewidth-32pt\relax}
\blue{\textbf{Key Idea}}\par\vspace{4pt}}
{\end{minipage}\egroup;\end{tikzpicture}\vspace{6pt}}

\newenvironment{redbox}
{\vspace{6pt}\noindent\begin{tikzpicture}
\node[
draw=myred, line width=1pt,
rounded corners=6pt,
inner sep=8pt,
text width=\dimexpr\linewidth-20pt\relax
]\bgroup
\begin{minipage}{\dimexpr\linewidth-32pt\relax}
\red{\textbf{Important}}\par\vspace{4pt}}
{\end{minipage}\egroup;\end{tikzpicture}\vspace{6pt}}

%%%%%%%%%%%%%%%%%%%%%%%%%%%%%%%%%%%%%%%%
% NOTE SPACE COMMAND
%%%%%%%%%%%%%%%%%%%%%%%%%%%%%%%%%%%%%%%%
\newcommand{\notespace}{
\vspace{0.8cm}
\hrule
\vspace{0.8cm}
}

%%%%%%%%%%%%%%%%%%%%%%%%%%%%%%%%%%%%%%%%
% SMALL HELPERS (for clean lists)
%%%%%%%%%%%%%%%%%%%%%%%%%%%%%%%%%%%%%%%%
\setlist[itemize]{leftmargin=1.3em}
\newcommand{\term}[1]{\blue{\textbf{#1}}}
\newcommand{\danger}[1]{\red{\textbf{#1}}}

%%%%%%%%%%%%%%%%%%%%%%%%%%%%%%%%%%%%%%%%
\begin{document}

\begin{center}
{\LARGE \blue{AI Infrastructure / GPU Systems Notes}}\\
\vspace{4pt}
{\small Live Lecture Notes Template}
\end{center}

%%%%%%%%%%%%%%%%%%%%%%%%%%%%%%%%%%%%%%%%
%%%%%%%%%%%%%%%%%%%%%%%%%%%%%%%%%%%%%%%%

\section{Module 2 — AI-Centric Data Center}

\subsection{Inside an AI Data Center}
There is compute so that whatever request comes in, we can process it. One server may not be enough, so we need multiple compute nodes and a network to communicate between them. Then we need storage for the data.

Support infrastructure includes power, cooling, security, etc.

\begin{bluebox}
\textbf{AI data center =} \blue{compute} + \blue{network} + \blue{storage} + \blue{support infra (power/cooling/security)}.
\end{bluebox}

\textbf{Key constraints when deploying high-density GPU workloads:}
\begin{itemize}
\item \term{Power:} limited power capacity per rack; GPUs need a high, consistent power supply.
\item \term{Heat:} cooling capacity of racks and rooms is a bottleneck.
\item \term{Space:} GPUs at scale require a large physical footprint (and good airflow/layout).
\end{itemize}
\notespace

\subsection{New Paradigms}
\begin{itemize}
\item \term{GPUs and parallelism:} scale performance by running many operations at once.
\item \term{Chiplets:} multiple smaller chips combined for scalability (yield/cost/perf benefits).
\item \term{3D stacking:} vertical layers to boost density (closer memory/compute).
\item \term{AI accelerators:} custom silicon for massive speedups on specific workloads.
\end{itemize}
\notespace

\subsection{Data Processing Unit (DPU)}
Handles \textbf{data-centric} tasks:

\textbf{Networking}
\begin{itemize}
\item packet processing
\item load balancing
\item overlay/underlay networking
\item RDMA (remote direct memory access)
\end{itemize}

\textbf{Storage}
\begin{itemize}
\item compression/decompression
\item encryption for data at rest
\item RAID/erasure coding
\item data dedupe
\end{itemize}

\textbf{Security}
\begin{itemize}
\item firewalling \& packet inspection
\item IPsec/TLS offload
\item multi-tenant isolation
\item zero trust enforcement
\end{itemize}

\begin{bluebox}
\textbf{Why DPUs exist:} offload data movement + infrastructure work from the CPU so the CPU/GPU can focus on application logic and training/inference.
\end{bluebox}

\subsubsection*{\blue{CPU vs GPU vs DPU (cleaned + clarified)}}
CPUs are best at \blue{decision-making and control flow}, but they don’t scale well to thousands of data-movement tasks.

GPUs are best for \blue{AI training and parallel computation}, but are not ideal for OS-level operations and general-purpose control logic.

DPUs take care of \blue{data movement and infrastructure offload}. They are best at the networking/storage/security tasks above, but not good for running heavy application compute.

\begin{center}
\renewcommand{\arraystretch}{1.15}
\begin{tabular}{|p{0.18\linewidth}|p{0.36\linewidth}|p{0.36\linewidth}|}
\hline
\textbf{Unit} & \textbf{Best at} & \textbf{Not great at} \\
\hline
CPU & control flow, OS, scheduling, general logic & massive parallel math throughput \\
\hline
GPU & matrix math, parallel compute, training/inference kernels & branchy logic, OS responsibilities \\
\hline
DPU & networking/storage/security offload, RDMA pipelines & heavy application compute \\
\hline
\end{tabular}
\end{center}
\notespace

\subsection{Networking}
We need separate networks for different tasks. Separation helps with \blue{performance isolation}, \blue{latency sensitivity}, \blue{security}, \blue{scaling}, and \blue{failure robustness}.
\notespace

\subsection{Network Fabric}
There is \textbf{logical} and \textbf{physical} networking.

\textbf{Compute fabric:}
\begin{itemize}
\item Designed primarily for GPU-to-GPU communication.
\item Implemented through InfiniBand, RoCE, or NVLink (depending on scope).
\item Requirements: extremely high throughput, low latency, and scalability.
\end{itemize}

\textbf{In-band management network (configuration/control within infra):}
\begin{itemize}
\item Handles control-plane management: SSH, job scheduling, access to code repos.
\item Typically Ethernet-based leaf/spine, with logical separation (VLANs).
\item Requirements: reliability, security/traffic isolation, moderate bandwidth.
\end{itemize}

\textbf{Out-of-band (OOB) management network:}
\begin{itemize}
\item Provides management even if the OS is down.
\item Uses a \term{Baseboard Management Controller (BMC)} to remotely manage the server (monitoring/logging/power cycling).
\item Requirements: always available, strong security/access control, separate physical ports, low-speed switches.
\end{itemize}

\begin{center}
\renewcommand{\arraystretch}{1.15}
\begin{tabular}{|p{0.22\linewidth}|p{0.34\linewidth}|p{0.34\linewidth}|}
\hline
\textbf{Network} & \textbf{Purpose} & \textbf{Key requirements} \\
\hline
Compute & GPU$\leftrightarrow$GPU traffic (training/inference) & \danger{low latency}, \danger{high throughput}, scalable \\
\hline
In-band mgmt & SSH, scheduling, repos, config & reliable, secure isolation, moderate BW \\
\hline
OOB mgmt & manage server if OS down (BMC) & \danger{always on}, strong access control \\
\hline
\end{tabular}
\end{center}

\notespace

\subsection{Ethernet vs InfiniBand}
\textbf{Ethernet:}
\begin{itemize}
\item General-purpose networking; uses TCP/IP stack.
\item Universally supported.
\item Typical speeds: 1--400 Gbps (depends on generation and deployment).
\end{itemize}

\textbf{InfiniBand:}
\begin{itemize}
\item Ultra-fast with specialized design for HPC.
\item Lower latency; supports RDMA (bypass CPU with low overhead).
\item Typical speeds: 10--400+ Gbps (depends on generation).
\item Requires specialized drivers; often more expensive.
\end{itemize}

\begin{bluebox}
\textbf{Simple mental model:} Ethernet = universal + flexible. InfiniBand = specialized + low-latency HPC.
\end{bluebox}
\notespace

\subsection{Converged Ethernet}
Converged LAN + SAN + HPC.
\begin{itemize}
\item Can use RDMA over Ethernet (e.g., RoCE) to bypass CPU.
\end{itemize}
\notespace

\subsection{Storage}
A data center uses \term{tiered storage}:

\begin{itemize}
\item \textbf{(1) NVMe SSD (local):} non-volatile local storage for fast I/O during training/inference; limited capacity but very fast.
\item \textbf{(2) Clustered storage / parallel file system:} shared high-speed access across many GPU nodes in a cluster.
\item \textbf{(3) Network file system (NFS):} distribute small datasets, configs, and scripts across nodes.
\item \textbf{(4) Object storage:} long-term storage for massive datasets, model checkpoints, and logs.
\end{itemize}

\begin{center}
\renewcommand{\arraystretch}{1.15}
\begin{tabular}{|p{0.20\linewidth}|p{0.42\linewidth}|p{0.30\linewidth}|}
\hline
\textbf{Tier} & \textbf{Used for} & \textbf{Tradeoff} \\
\hline
Local NVMe & hot data + fast scratch I/O & fastest, limited capacity \\
\hline
Parallel FS & shared high-speed access & scalable, more complex \\
\hline
NFS & configs/scripts/small data & simple, not for heavy throughput \\
\hline
Object store & checkpoints/logs/big datasets & durable + cheap, higher latency \\
\hline
\end{tabular}
\end{center}
\notespace

\subsection{Cloud vs On-Prem}
\textbf{Cloud:}
\begin{itemize}
\item Accessibility: no need to buy infrastructure.
\item Pay-as-you-go.
\item Flexible scaling.
\item Control depends on the provider.
\end{itemize}

\textbf{On-prem:}
\begin{itemize}
\item Data security and sovereignty.
\item High upfront cost.
\item Limited by owned hardware.
\item Full control over data and systems.
\end{itemize}
\notespace

%%%%%%%%%%%%%%%%%%%%%%%%%%%%%%%%%%%%%%%%
%%%%%%%%%%%%%%%%%%%%%%%%%%%%%%%%%%%%%%%%
\section{Module 3 — NVIDIA Technology Stack}

\subsection{NVIDIA AI Leadership}
\notespace

\subsection{Technology Stack Overview}
\notespace

\subsection{Layer 1 — Physical Infrastructure}
Here are the layers in NVIDIA infrastructure:
\begin{enumerate}
\item \term{Physical layer:} GPU, data center, NICs, DPU
\item \term{Data movement \& I/O acceleration:} NVLink, RDMA, storage, HPC fabrics (InfiniBand)
\item \term{OS / drivers / virtualization:} GPU drivers
\item \term{Core libraries:} CUDA, NCCL
\item \term{Monitoring \& management:} NVIDIA-SMI, etc.
\item \term{Applications / vertical solutions:} NVIDIA NIMs, enterprise stacks
\end{enumerate}
\notespace

\subsubsection{SuperPOD}
An architecture implementation of multiple DGX nodes: a scalable AI system.

\textbf{Rough flow:} jump box $\rightarrow$ management nodes $\rightarrow$ InfiniBand compute fabric $\rightarrow$ compute nodes $\rightarrow$ InfiniBand storage fabric $\rightarrow$ high-speed storage
\notespace

\subsubsection{ConnectX Networking}
Think of NICs as high-performance interconnect adapters for InfiniBand/Ethernet, used for reliable high-speed networking.
\notespace

\subsection{GPU Architecture Details}
\notespace

\subsubsection{GPU Cores}
\begin{itemize}
\item \term{CUDA cores:} general compute + graphics rendering; flexible for general parallel tasks.
\item \term{Tensor cores:} deep learning acceleration; major throughput for training/inference (often large speedups for dense math).
\item \term{Ray tracing cores:} real-time ray tracing/light simulation (graphics realism: shadows/reflections).
\end{itemize}

\begin{center}
\renewcommand{\arraystretch}{1.15}
\begin{tabular}{|p{0.22\linewidth}|p{0.52\linewidth}|p{0.18\linewidth}|}
\hline
\textbf{Core type} & \textbf{What it does} & \textbf{Where used} \\
\hline
CUDA core & general-purpose parallel compute & HPC, graphics, kernels \\
\hline
Tensor core & dense matrix math acceleration & training, inference \\
\hline
RT core & ray/triangle intersection + lighting & graphics \\
\hline
\end{tabular}
\end{center}

\notespace

\subsection{Layer 2 — Data Movement / IO}
\notespace

\subsubsection{NVLink}
Connection between multiple GPUs or (in some systems) CPU$\leftrightarrow$GPU. PCIe is slower; NVLink is a higher-bandwidth interconnect with both hardware and software components.

\begin{itemize}
\item NVLink bridges: consumer and professional graphics cards (some generations).
\item Integrated NVLink: NVLink inside the silicon.
\item NVSwitch: for large-scale multi-GPU systems.
\end{itemize}
\notespace

\subsubsection{InfiniBand}
High-performance connection used for distributed GPU clusters (fast inter-node communication). Often paired with RDMA for low overhead and high throughput.
\notespace

\subsubsection{RDMA}
Start with DMA: a hardware feature that lets supported devices transfer data to/from memory without the CPU acting as the middleman.

\begin{bluebox}
\textbf{DMA:} device $\leftrightarrow$ local host memory without CPU copies.\\
\textbf{RDMA:} device $\leftrightarrow$ \emph{remote} host memory (across the network) with low CPU overhead.
\end{bluebox}

RDMA is an extension of DMA: devices can bypass the CPU and access not only local host memory but also memory on another host.

Example mental model: host1 GPU $\rightarrow$ host2 system memory (via RDMA).
\notespace

\subsubsection{GPUDirect}
NVIDIA uses RDMA to implement GPUDirect features.

\textbf{GPUDirect RDMA:}
\begin{itemize}
\item Across hosts: host1 GPU memory $\rightarrow$ RDMA NIC $\rightarrow$ network $\rightarrow$ RDMA NIC $\rightarrow$ host2 GPU memory.
\item Bypasses OS, system memory, and CPU copies.
\item Ultra-low-latency networking between GPU nodes.
\item Needs an NVIDIA GPU + RDMA-capable NIC.
\end{itemize}

\textbf{GPUDirect Storage:}
\begin{itemize}
\item Within a host (and sometimes via storage fabric): GPU $\leftrightarrow$ storage, bypassing CPU bottlenecks.
\item High-throughput loading of large datasets into GPUs (e.g., NVMe RAID, parallel storage).
\item Needs an NVIDIA GPU + NVMe/parallel storage setup.
\end{itemize}

\begin{center}
\renewcommand{\arraystretch}{1.15}
\begin{tabular}{|p{0.26\linewidth}|p{0.46\linewidth}|p{0.20\linewidth}|}
\hline
\textbf{Feature} & \textbf{Data path (simplified)} & \textbf{Goal} \\
\hline
GPUDirect RDMA & GPU mem $\leftrightarrow$ NIC $\leftrightarrow$ network $\leftrightarrow$ NIC $\leftrightarrow$ GPU mem & low latency, avoid CPU copies \\
\hline
GPUDirect Storage & GPU mem $\leftrightarrow$ NVMe / storage fabric & feed GPUs faster (I/O) \\
\hline
\end{tabular}
\end{center}

\notespace

%%%%%%%%%%%%%%%%%%%%%%%%%%%%%%%%%%%%%%%%
\subsection{Layer 3 — OS / Drivers / Virtualization}
%%%%%%%%%%%%%%%%%%%%%%%%%%%%%%%%%%%%%%%%
\notespace

\subsubsection{GPU Drivers}
Software that connects NVIDIA GPUs with the operating system and applications.
\begin{itemize}
\item Enables GPU acceleration for AI/HPC/graphics.
\item Manages GPU resources, memory, and performance.
\item Supports containers, virtualization, and multi-GPU setups.
\end{itemize}
\notespace

\subsubsection{Virtualization}
If a physical server has a GPU and CPU, we can add a virtualization layer and create a virtual CPU/VM and a virtual GPU.

Virtualization matters because multiple users/workloads can share the same hardware efficiently. It enables isolation, flexible management, better utilization, and reproducibility.

\textbf{CPU virtualization:} divides CPU resources into multiple VMs, each running its own OS and applications with minimal interference.

\textbf{GPU virtualization:} more complex; requires techniques like GPU pass-through, vGPU, or API interception to safely share a GPU across workloads (common in cloud environments).
\notespace

\subsubsection{MIG vs vGPU}
\textbf{vGPU (software virtualization):}
\begin{itemize}
\item Physical GPU $\rightarrow$ hypervisor (NVIDIA virtualization software) $\rightarrow$ virtual GPUs (up to many) $\rightarrow$ VMs.
\item Hypervisor-based (e.g., VMware); uses guest drivers + a vGPU manager.
\end{itemize}

\textbf{MIG (Multi-Instance GPU):}
\begin{itemize}
\item Physical GPU $\rightarrow$ Linux + NVIDIA tools (e.g., \texttt{nvidia-smi}) + container runtime (Docker) $\rightarrow$ MIG instances (e.g., up to 7 on some GPUs) $\rightarrow$ containers.
\item Hardware-enforced strong isolation; common for data centers, HPC, AI/ML containers.
\end{itemize}

\begin{center}
\renewcommand{\arraystretch}{1.15}
\begin{tabular}{|p{0.22\linewidth}|p{0.34\linewidth}|p{0.34\linewidth}|}
\hline
\textbf{Tech} & \textbf{Isolation} & \textbf{Typical usage} \\
\hline
vGPU & software-enforced (hypervisor) & VM environments, cloud multi-tenant VMs \\
\hline
MIG & hardware-enforced & containers, HPC/AI clusters, strong partitioning \\
\hline
\end{tabular}
\end{center}

\notespace

%%%%%%%%%%%%%%%%%%%%%%%%%%%%%%%%%%%%%%%%
\subsection{Layer 4 — Libraries}
%%%%%%%%%%%%%%%%%%%%%%%%%%%%%%%%%%%%%%%%
\notespace

\subsubsection{cuDNN / NCCL}
\textbf{NCCL:} multi-GPU communication library that provides abstractions and optimizes patterns of communication between many GPUs.

Think of it like a traffic management system: it coordinates transfers and collective operations (e.g., gradient synchronization) across many GPUs.

\notespace

%%%%%%%%%%%%%%%%%%%%%%%%%%%%%%%%%%%%%%%%
\subsection{Layer 5 — Monitoring}
%%%%%%%%%%%%%%%%%%%%%%%%%%%%%%%%%%%%%%%%
\notespace

\subsubsection{NVIDIA-SMI}
Used for quick GPU monitoring and basic management on a single node. It reports GPU utilization and memory usage.
\notespace

\subsubsection{BCM (Base Command Manager)}
Manage, monitor, and orchestrate GPU resources and workloads across an entire AI data center:
\begin{itemize}
\item Cluster resource utilization, infra health, job status, user quotas, workload performance.
\item Management: firmware configs, power policies, software updates, provisioning, scheduling, allocation, deployments.
\item Interfaces: CLI, web UI, REST API.
\end{itemize}

Installation/deployment is typically an enterprise multi-component system requiring dedicated management nodes.

\subsubsection*{\blue{Monitoring capabilities (cleaned)}}
\begin{itemize}
\item CPU, NVIDIA GPUs, NVIDIA networking, storage (NFS, clusters, NVMe), workloads (AI/HPC).
\end{itemize}

\subsubsection*{\blue{Management/config capabilities (cleaned)}}
\begin{itemize}
\item GPU firmware configs, power policies, allocation.
\item Workload submission/orchestration through Slurm/Kubernetes.
\end{itemize}

\notespace


%%%%%%%%%%%%%%%%%%%%%%%%%%%%%%%%%%%%%%%%
\end{document}
